% !TeX spellcheck = de_DE 
\addscenariosection{1}{Kooperatives Szenario}{Angriff aus der Tiefe}{\images/angriff_aus_dem_untergrund.png} 

\begin{multicols*}{2} 

\textbf{Autor:} dwarftoe 

\textit{Die Erde bebt! Oh nein ... jetzt hat es auch hier begonnen. Flüchtlinge aus einem fernen Königreich hatten von schrecklichen Erdbeben berichtet, die ihre gesamte Heimat verwüstet hatten. Ganze Landstriche waren im bebenden Boden versunken. Berge waren aufgebrochen, Wälder eingestürzt und die Erde selbst war auseinandergerissen worden. Dann, ohne jede Vorwarnung, tat sich vor den Städten der Boden auf. Aus den Tiefen unter der Welt strömten Horden monströser Kreaturen hervor. Wie eine lebendige Flut fegten sie über das Land, setzten alles in Brand und stürmten die Städte. Keine Armee war auf einen solchen Feind vorbereitet gewesen. Keine Mauer hatte ihnen standhalten können. Königreich um Königreich versank in Trümmern. Nur eine Handvoll Menschen überlebte die Verwüstung. Sie flohen so weit sie konnten und trugen die schrecklichen Geschichten dessen, was sie erlebt hatten, in die Welt hinaus. Doch nun ... geschieht dasselbe hier. Die Erde bebt unter den Füßen der Menschen. Risse breiten sich über den Boden aus. Tief unter der Oberfläche regt sich etwas. Den Königreichen bleibt nur wenig Zeit zum Handeln. Zum ersten Mal legen ihre Herrscher ihre uralten Feindschaften beiseite und schließen ein verzweifeltes Bündnis. Ihre Armeen sammeln sich, die Tore werden verriegelt und ihre Krieger bereiten sich darauf vor, dem unbekannten Grauen entgegenzutreten, das aus den Tiefen der Erde emporsteigt. Doch es gibt eine schreckliche Frage, auf die niemand eine Antwort weiß: Wie bekämpft man einen Feind, dessen Angriff man nicht kommen sieht? Erneut bebt der Boden. Etwas kommt. Etwas Uraltes. Etwas Hungriges. Macht euch bereit für das Ende.} 

\subsection*{\MakeUppercase{Szenariodauer}} 5 + $2P$ Runden ($P$ steht für die Anzahl der Spieler)

\subsection*{\MakeUppercase{Spieleraufstellung}} 
\textbf{Spieleranzahl:} 1--5 

\textbf{Startressourcen:} 13 \svg{gold}, 4 \svg{building_materials}, 1 \svg{valuables} 

\textbf{Start-Einkommen:} 10 \svg{gold}, 0 \svg{building_materials}, 0 \svg{valuables} 

\textbf{Starteinheiten:} 
\begin{itemize} 
\item Ein ``Rudel'' der günstigsten \svg{bronze}-Einheiten 
\item "Wenige" der teuersten \svg{bronze}-Einheiten 
\end{itemize} 

\textbf{Stadtgebäude:} \svg{bronze}~Gebäude, Rathaus, \svg[12]{building_special_gears}~Gebäude 

\textbf{Kartenteile-Vorrat:} Jeder Spieler erhält 1 zufälliges fernes Gebiet (II--III) 

\subsection*{\MakeUppercase{Kartenaufbau}} Nehmt die folgenden Gebiete und ordnet sie wie in der Szenariokarte dargestellt an ($P$ steht für die Anzahl der Spieler): 

\begin{itemize} 
\item $P$ × Startgebiet (I) 
\item 1 × Startgebiet (I) Dungeon 
\item $2P$ × ferne Gebiete (II--III) 
\item $2P$ × nahe Untergrundgebiete (IV--V) 
\item 1--3 × zentrale Untergrundgebiete (VI--VII) (abhängig von $P$) 
\item 1--2 × nahe Untergrundgebiete (IV--V) (abhängig von $P$) 
\end{itemize} 

\subsection*{\MakeUppercase{Siegbedingungen}} Ihr müsst in den Untergrund vordringen und den bösen Kriegsmagier in seiner Stadt besiegen. \subsection*{\MakeUppercase{Niederlagebedingungen}} Am Ende der Runde 5 + $2P$ haben die Spieler verloren. Wird eine Stadt von der feindlichen KI-Armee eingenommen, haben die Spieler 2 Runden Zeit, um sie zurückzuerobern, ansonsten haben auch alle verloren. Ist eine Stadt eingenommen und der zugehörige Haupt-Held auch besiegt, endet das Spiel sofort.
\subsection*{\MakeUppercase{Zeitgesteuerte Ereignisse}} 
\textbf{4. Runde:} 
\begin{itemize} \item Ein zufälliges Gebiet II--III, das an ein Untergrundgebiet angrenzt, wird durch ein Untergrundgebiet IV--V ersetzt. Um das Gebiet auszuwählen, könnte man einen kleinen Würfel auf die Karte fallen lassen und das Gebiet, das am nächsten zum Würfel liegt, wird ausgetauscht. Oder man nummeriert die Gebiete und würfelt dann mit einem W4, W6, W8 oder W10, je nach Anzahl der Gebiete und tauscht dieses aus.
\end{itemize} 
\textbf{In jeder darauffolgenden Runde, bis alle Gebiete II--III, die an ein ursprüngliches Untergrundgebiet angrenzten, ersetzt wurden:} 
\begin{itemize} 
\item Ein zufälliges Gebiet II--III, das an ein Untergrundgebiet angrenzt, wird durch ein Untergrundgebiet IV--V ersetzt. 
\end{itemize} 
\subsection*{\MakeUppercase{Zusätzliche Regeln}} 
\begin{itemize} 
\item Das Startgebiet (I) Dungeon zählt als Untergrundgebiet und kann erst betreten werden, wenn sich alle Untergrundtore geöffnet haben.
\item Alle Spieler dürfen einmal pro Runde oder wenn sie auf einem Handelsposten stehen, mit einem weiteren Spieler Gold und Resourcen tauschen.
\item Alle Spieler dürfen Einheiten, Artefakte und Zauber tauschen, wenn die Helden sich auf angrenzenden Feldern befinden.
\item Wird ein Gebiet durch ein Untergrundgebiet ersetzt, werden alle Würfel entfernt. Das Einkommen wird entsprechend angepasst. Das Untergrundgebiet gilt als unerforscht, es sei denn, ein Held befindet sich auf diesem Gebiet. 
\item Befindet sich ein Held auf dem ausgetauschten Gebiet, stürzt er in den Untergrund. Das Gebiet gilt als erforscht, aber die schwächste Einheit erleidet Schaden: Rudel \rightarrow Wenige, bzw. Wenige \rightarrow tot. 
\item Solange es keinen Weg zurück zur Heimatstadt gibt, kann euer unterirdischer Held keine Einheiten aus seiner Heimatstadt rekrutieren. Gewinnt der Held jedoch einen Kampf gegen neutrale Einheite, kann er eine der besiegten Einheiten für die Hälfte der Kosten in seine Armee aufnehmen. 
\item Solange es keinen Weg von eurem unterirdischen Helden zurück zur Heimatstadt gibt, könnt ihr einen zweiten Helden für die Hälfte der Kosten rekrutieren. Alle Einheiten, die ihr in eurer Stadt rekrutiert, bleiben entweder in der Stadt oder schließen sich eurem Nebenhelden an. 
\item Der Nebenheld ist Stufe I und erhält die Starteinheiten + die Einheiten, die in der Stadt rekrutiert wurden (nur die, die der Hauptheld nicht hat)
\item Wenn das zweite Untergrundgebiet an euer Startgebiet angrenzt, öffnet sich darauf ein Untergrundtor, das an eure Stadt angrenzt. Darauf erscheint eine feindliche Armee. Gleich am Anfang der nächsten Runde greift die feindliche Armee eure Stadt an. Man kann die Armee auch gleich in der Runde, in der sie auftaucht, am Untergrundtor angreifen.
\item Die Stärke der feindlichen Armee hängt von der Runde ab, in der sie angreift/angegriffen wird. Die Armee besteht aus neutralen Einheiten der Dungeon-Partei. Alle zwei Runden erhöht sich die Anzahl der Einheiten um eine Einheit. Zuerst kommen zwei \svg{bronze}-Einheiten (Harpien und Augen des Bösen), danach zwei \svg{silver}-Einheiten (Medusen und Minotauren) und anschließend zwei \svg{golden}-Einheiten (Mantikore und Schwarze Drachen) hinzu. Für den angreifenden Helden werden alle zwei Runden eine zufällige KI-Karte mit Angriff oder Verteidigung hinzugefügt. (Beispiel: In Runde 10 besteht die Armee aus den 2 \svg{bronze}-Einheiten, den 2 \svg{silver}-Einheiten und 1 \svg{golden}-Einheit (Mantikor). Der KI-Stapel umfasst 5 zufällige Angriffs- oder Verteidigungskarten.) 
\item Ist euer Hauptheld nicht in der Stadt, könnt ihr 10 Gold bezahlen, um eure Einheiten in eure Stadt zu versetzen. Stürzte euer Hauptheld in den Untergrund, könnt ihr alternativ eure Stadt mit den dort rekrutierten Einheiten (+ evtl. Nebenheld) verteidigen. Wurde die Stadt verteidigt, ist der Weg frei und alle verbliebenen Einheiten können zum Hauptheld transferiert werden. Als Belohnung erhält der Spieler 5 \svg{gold}, 1 \svg{artifact} und \svg{morale_positive}.
\item Verliert ihr die Stadt, generiert sie kein Einkommen mehr und man kann keine Einheiten mehr rekrutieren. Sie muss dann innerhalb von 2 Runden von einem Spieler zurückerobert werden, sonst ist das Szenario verloren.
\item Nachdem sich das letzte Untergrundtor geöffnet hat, bleiben den Spielern noch zwei Runden, um die feindliche Dungeon-Stadt zu besiegen. Andernfalls verlieren sie. 
\item Die letzte Schlacht findet in der feindlichen Dungeon-Stadt (mit Mauern und Geschützturm) statt. Die Stärke der feindlichen Armee hängt wieder von der Rundenanzahl ab. Die Armee besteht aus den Dungeon-Einheiten und wird folgendermaßen festgelegt: 

Runde 1 -- Wenige \svg{bronze} Troglodyten, 

Runde 2 -- Rudel \svg{bronze} Troglodyten, 

Runde 3 -- Rudel \svg{bronze} Troglodyten + Wenige \svg{bronze} Harpien, 

Runde 4 -- Rudel \svg{bronze} Troglodyten + Rudel \svg{bronze} Harpien, usw.

\hspace{1em}\vdots

Runde 14 -- Alle Dungeon-Einheiten auf der Rudel-Seite.

Um zu gewinnen, müssen alle Einheiten besiegt werden. Sind es mehr als 5 Einheiten, werden 5 zufällig ausgewählt und aufgestellt. Wird einen davon besiegt, wird eine weitere auf das Spielfeld gestellt. Der feindliche Held bekommt entsprechend der Rundenanzahl eine gewisse Anzahl KI-Bot-Karten, nämlich nach folgendem Muster: Für Runde 1 und 2 je eine Zauberkarte, für Runde 3 eine Fähigkeitenkarte, für Runde 4 eine Angriffskarte, für Runde 5 eine Verteidigungskarte. Dann wieder von vorne, also Runde 6 und 7 eine Zauberkarte, Runde 8 Fähigkeit, dann Angriff, Verteidigung, 2 Zauber, Fähigkeit, usw. Die Fähigkeit ist Zauberei, d.h. der nächste Zauber erhält +2 \svg{empower}. Die Zauber sind zufällig aus der Erdmagieschule (Steinhaut, Implosion, Sorge...) ausgewählt. Verliert ein Spieler die Schlacht, bleiben die besiegten Einheiten besiegt und ein nächster Spieler kann die Schlacht fortsetzen. Der KI-Stapel und die Zauber werden neu gemischt.
\end{itemize} 
\end{multicols*} 
\vspace{3em} 
\begin{center} 
\includegraphics[width=0.5\paperwidth]{\maps/angriff_aus_dem_untergrund_5p.png} \captionof{figure}{\textbf{SZENARIO FÜR 5 SPIELER}} \end{center} 
\vspace{3em} 
\begin{center} 
\includegraphics[width=0.5\paperwidth]{\maps/angriff_aus_dem_untergrund_4p.png} \captionof{figure}{\textbf{SZENARIO FÜR 4 SPIELER}} \end{center} 
\vspace{3em} 
\begin{center} 
\includegraphics[width=0.38\paperwidth]{\maps/angriff_aus_dem_untergrund_3p.png} \captionof{figure}{\textbf{SZENARIO FÜR 3 SPIELER}} \end{center} 
\vspace{1em} 
\begin{center} 
\includegraphics[width=0.28\paperwidth]{\maps/angriff_aus_dem_untergrund_2p.png} \captionof{figure}{\textbf{SZENARIO FÜR 2 SPIELER}} \end{center} 
\vspace{1em} 
\begin{center} 
\includegraphics[width=0.2\paperwidth]{\maps/angriff_aus_dem_untergrund_1p.png} \captionof{figure}{\textbf{SZENARIO FÜR 1 SPIELER}} \end{center}